\documentclass[12pt]{article}
\usepackage[utf8]{inputenc}
\usepackage[T1]{fontenc}
\usepackage{hyperref}
\title{Quebec Sign Language Language Resources for K-12 Curriculum Development}
\author{Compiled by Claude Sonnet 5.5 (claude-sonnet-5-5) for Andr\'as Kornai}
\date{}
\begin{document}
\maketitle

\section{Introduction}
Quebec Sign Language (Langue des signes qu\'eb\'ecoise, LSQ; ISO 639-3 fcs) is the sign language of the Deaf community of Quebec. This dossier is based on a small number of web searches, which turned up mostly policy and academic reports about the status and teaching of LSQ. No LSQ dictionary, textbook or school curriculum was found in the search results, and none is cited here.

\section{Primary Educational Resources}
No primers, textbooks or classroom materials for LSQ were found in the search results. A report from a Quebec working group on the teaching of LSQ (Office des personnes handicap\'ees du Qu\'ebec) bears on its use as a language of instruction:
\begin{itemize}
\item OPHQ report on LSQ: \url{https://www.ophq.gouv.qc.ca/fileadmin/centre_documentaire/Etudes__analyses_et_rapports/OPHQ_Rapport_LSQ.pdf}
\end{itemize}

\section{Linguistic and Reference Materials}
No dictionary or grammar of LSQ appeared in the results.

\section{Religious and Community Resources}
No community or religious resources were verified in the search results.

\section{Digital and Audio Resources}
No video or digital LSQ resources were found in the search results.

\section{Curriculum Development Context}
The following documents discuss the recognition, history and educational use of LSQ:
\begin{itemize}
\item OPHQ report on the recognition of sign languages: \url{https://www.ophq.gouv.qc.ca/fileadmin/centre_documentaire/Etudes__analyses_et_rapports/RAP_reconnaissance_langues_signes.pdf}
\item Article in the journal \emph{Globe} (2013) on the difficulty of having LSQ recognised as a first language of learning: \url{https://www.erudit.org/en/journals/globe/2013-v16-n1-globe0796/1018177ar.pdf}
\item Luna (2015), on methods of teaching deaf pupils in Quebec, 1850--1960: \url{https://cla-acl.artsci.utoronto.ca/pdfs/actes-2015/Luna-2015.pdf}
\end{itemize}
A K-12 curriculum would need to start from primary materials that these searches did not locate; targeted searches of Deaf community organisations and universities in Quebec are recommended.

\end{document}